\section{Continuum and infinite-volume limits}
\label{sec:continuum}

The source estimates provide convergence on descendants of fixed coarse
cells.  We first obtain bounds that allow these cells to approximate
arbitrary test functions.  The trace estimates then identify the
infinite-volume limit, and comparison of the two inclined blockings
supplies the rotations absent from the microscopic lattice.

Write $a_N=L^{-N}$ and let $B_N>0$ be the field normalization in
Equation~\eqref{eq:field-normalization}.  For a compactly supported
function $f:\R^2\to\R^3$, define
\begin{equation}
 \phi_N(f)=B_Na_N^2\sum_{x\in\Lambda_N}
                  f(x)\cdot q_x,
 \label{eq:smeared-lattice-field}
\end{equation}
where $\Lambda_N$ is the physical lattice, with its specified origin and
orientation.  On a torus the sum uses its periodic lattice and a periodic
test function.  Scalar tests of a specified component have the same
meaning.  For component indices
$\boldsymbol\alpha=(\alpha_1,\ldots,\alpha_n)\in\{1,2,3\}^n$, introduce
the lattice moment measure
\begin{equation}
 S_{n,N}^{\boldsymbol\alpha}
   =(B_Na_N^2)^n\sum_{x_1,\ldots,x_n\in\Lambda_N}
       \E\!\left[\prod_{i=1}^n q_{x_i}^{\alpha_i}\right]
       \delta_{(x_1,\ldots,x_n)}.
 \label{eq:lattice-moment-measure}
\end{equation}
The expectation will be taken either on one of the physical tori of
Section~\ref{sec:trace} or in the periodic infinite-plane state at the
same cutoff.  These choices will always be indicated when they matter.

\subsection{Local moment bounds}

A terminal cell is the set of microscopic descendants of one layer-zero
site in a standard blocking.  In physical coordinates it is a unit
square, up to the choice of lattice boundary convention.  We use the
term unit box for any translate of a unit square in the regulator axes.
On a torus, boxes are interpreted in a periodic coordinate chart;
bounded-multiplicity coverings give the same estimates when a chart
crosses a period.

\begin{proposition}[Moment majorants]
\label{prop:moment-majorants}
For every $n\ge1$ and every component array $\boldsymbol\alpha$, the
measure $S_{n,N}^{\boldsymbol\alpha}$ is nonnegative.  There is a constant
$C$, independent of $N\ge K_0$, the axis torus in the exhaustion of
Section~\ref{sec:trace}, and the positions of the unit boxes, such that
\begin{equation}
 S_{n,N}^{\boldsymbol\alpha}(Q_1\times\cdots\times Q_n)
       \le C^n n!
 \label{eq:unit-box-majorant}
\end{equation}
for all unit boxes $Q_1,\ldots,Q_n$.  The same bound holds in the
periodic infinite-plane state.

For each fixed number $D$ of terminal cells, their component field sums
have joint exponential moments in a neighborhood of the origin,
uniformly in these cutoffs and axis tori.  If $\mu$ is the source-free
terminal probability law and $z_Y\in\C^3$ is supported on at most $D$
terminal sites with $|z_Y^\alpha|\le1$, their generating function differs
from
\begin{equation}
             \E_\mu\exp\!\left(\sum_Y z_Y\cdot V_Y\right)
 \label{eq:orig-34}
\end{equation}
by $o_H(1)$, uniformly for fixed $D$.

On each fixed tilted physical torus of Section~\ref{sec:trace}, the
exponential-moment bounds and Equation~\eqref{eq:unit-box-majorant}
hold with constants that may depend on that torus and on $H$, but remain
independent of $N$.
\end{proposition}

\begin{proof}
For finite volume, expand each interaction factor as
\[
 e^{\beta_N q_x\cdot q_y}
   =\prod_{\alpha=1}^3\sum_{m\ge0}
       \frac{\beta_N^m}{m!}(q_x^\alpha q_y^\alpha)^m.
\]
The expansion is absolutely convergent.  After multiplication by any
specified spin components, each single-site integral is zero if one
coordinate has odd degree, and is nonnegative otherwise.  Since
$\beta_N\ge0$, every surviving term is nonnegative.  This proves the
first assertion, including repeated insertion sites and oblique
periods.  Positivity passes to the fixed-cutoff periodic plane limit.

Let $J$ be a set of at most $D$ terminal sites, and let $C_Y$ be the
microscopic descendants of $Y\in J$.  Iterating the exact source
substitution of Proposition~\ref{prop:source-step} identifies $z_Y$ with
the coupling to the vector
\[
       B_Na_N^2\sum_{x\in C_Y}q_x.
\]
At the terminal layer, write the source-dependent density, after its
common bulk scalar has been removed, as
\[
 e^{X(V)}
 e^{\sum_{Y\in J}z_Y\cdot V_Y+G(V;z)}\Xi_z(V),
\]
where $e^X\Xi_0$ has normalizer $Z$ and law $\mu$.  The source-anchor
convention and the norm bound of Proposition~\ref{prop:source-step}
give
\[
       \sup_V|G(V;z)|\le DR_0.
\]
Indeed, a monomial that survives setting all sources outside $J$ to
zero has its designated source anchor in $J$.  The same specialization
has a useful consequence for the gas: every surviving changed label
has complete support meeting $J$.  Its weighted activity difference
therefore satisfies
\[
 \sum_\lambda e^{2s_\lambda}
            \norm{k_\lambda(z)-k_\lambda(0)}_\infty
       \le Dw_0.
\]
The original and source-dependent activities satisfy the combined
site-hit hypothesis of Equation~\eqref{eq:orig-16}.  Consequently
\[
 \frac1Z\int e^X|\Xi_z-\Xi_0|
      \le e^{Dw_0}-1.
\]
Since $|\sum_{J}z_Y\cdot V_Y|\le3D$, subtracting
Equation~\eqref{eq:orig-34} from the normalized source integral gives
\begin{equation}
 \left|\E e^{\sum_{Y\in J}z_Y\cdot\phi_N(C_Y)}
       -\E_\mu e^{\sum_{Y\in J}z_Y\cdot V_Y}\right|
 \le e^{3D}(e^{DR_0}-1)
      +e^{3D+DR_0}(e^{Dw_0}-1),
 \label{eq:localized-source-comparison}
\end{equation}
where $\phi_N(C_Y)=B_Na_N^2\sum_{x\in C_Y}q_x$ is vector-valued.
The right side is $o_H(1)$ for fixed $D$.  In particular the generating
functions are bounded on a fixed complex polydisc, independently of
volume and cutoff.  This is a localized use of the integrated gas
comparison: its bound depends on $D$, rather than on the torus volume.

For one real component cell sum $X$, the bounds for $\E e^X$ and
$\E e^{-X}$ imply $\E e^{|X|}\le C_0$.  Thus
$\E|X|^n\le C_0n!$.  If $X_i$ are component sums in possibly different
terminal cells, H\"older's inequality yields
\[
  \left|\E\prod_{i=1}^n X_i\right|
      \le\prod_{i=1}^n(\E|X_i|^n)^{1/n}
      \le C^n n!.
\]
The left side without absolute values is the mass of the corresponding
product of cells.  An arbitrary unit box is covered by a fixed number
of terminal cells.  Nonnegativity of the moment measure therefore
extends the estimate to arbitrary translated products of unit boxes,
with a fixed enlargement of $C$.  This also shows that the estimates
are unchanged by translating the microscopic origin.  At a fixed
cutoff, the cell sums are bounded local observables, so the estimates
pass to the periodic plane state.

For a fixed tilted torus, use standard blocking in the regulator axes.
Its period lattice is divisible at every step, and its terminal volume
is fixed.  Absolute bounds on the regular exponent and on the covering
activities from Sections~\ref{sec:rg} and~\ref{sec:sources} bound the
source numerator in that volume.  The small-cap lower bound for the
source-free normalizer bounds the denominator away from zero after
removal of the bulk scalar.  This argument uses no reflection
positivity of the tilted torus.  It gives a finite bound depending on
its fixed terminal volume and on $H$, uniformly in depth.  The preceding
exponential-moment and covering arguments now apply there as well.
\end{proof}

We record a direct consequence that will be used repeatedly.  If real
component tests $f_1,\ldots,f_r$ have support in a fixed bounded region,
then every real linear combination of their smeared fields has moments
bounded by $C(f)^n n!$, uniformly over the cutoffs and axis tori under
consideration.  To see this, first bound an even moment by integrating
the absolute test coefficients against the nonnegative component
measures and applying Equation~\eqref{eq:unit-box-majorant}.  Odd
absolute moments follow by Cauchy--Schwarz from even ones, with an
exponential change of the constant.  These bounds imply exponential
tails at a positive radius.  On a fixed tilted torus the same conclusion
holds with a torus-dependent constant.

\subsection{Convergence and comparison of orientations}

\begin{proposition}[Continuum limit]
\label{prop:continuum-limit}
On every fixed axis or tilted physical torus used in
Section~\ref{sec:trace}, all joint moments of the fields
$\phi_N(f)$ converge for continuous tests.  Their finite-dimensional
laws converge, and the limiting laws are determined by their moments.
The moment measures converge on arbitrary continuous compactly
supported tests in the insertion coordinates.

In the periodic plane state these assertions hold for compactly
supported tests.  The same plane limit is obtained by taking the
continuum limit on the axis tori and then their exhaustion, or by taking
any simultaneous sequence of cutoffs and axis tori for which both the
cutoff depth and the doubling index tend to infinity.

Finally, let $O_+$ be the rotation with cosine $3/5$ and sine $4/5$.
The plane limits obtained with regulator orientations $\Id$ and
$O_+^2$ have identical finite-dimensional laws when tested in the same
physical coordinates.
\end{proposition}

\begin{proof}
\emph{Fixed tori and fixed coarse cells.}
Fix a physical torus and a layer $j$.  A test that is constant on the
microscopic descendants of its layer-$j$ cells has an exact source
representation at that layer.  The coefficient multiplying each test
value in the layer-$j$ source is
\begin{equation}
             L^{-2j}\prod_{i=1}^j c_{i,N}^{-1}.
 \label{eq:regular-cell-source-multiplier}
\end{equation}
Indeed, its microscopic coefficient is $B_NL^{-2N}$, while applying the
source substitution through the first $N-j$ steps multiplies this by
$L^{2(N-j)}\prod_{i=j+1}^N c_{i,N}$.  The definition
$B_N=\prod_{i=1}^Nc_{i,N}^{-1}$ gives
Equation~\eqref{eq:regular-cell-source-multiplier}.

Corollary~\ref{cor:source-endpoints} and the density comparison
Equation~\eqref{eq:orig-14} imply convergence of the normalized source
integrals near the source origin.  The multiplier in
Equation~\eqref{eq:regular-cell-source-multiplier} also converges,
because it contains a fixed number of factors.  To justify taking
ratios here, observe that at layer $j$ the torus has fixed volume.
Absolute gas bounds and regular-exponent bounds control the numerator,
and the small-cap lower bound controls the source-free denominator,
uniformly in the original depth $N$.  The common bulk scalar cancels
from numerator and denominator.  Thus source generating functions
converge uniformly on a sufficiently small complex polydisc.  Cauchy's
formula gives convergence of every joint moment of these cell tests.

The layer-$j$ grids can be aligned across depths.  One may choose the
physical origins of the approximating lattices, or translate successive
block partitions: their available shifts fill one ancestral block
modulo the microscopic spacing.  The period identifications are
compatible with these choices.  An origin error of size $O_L(a_N)$
has vanishing effect on any smooth test by the moment majorants and
uniform continuity.  More generally, the grids used in a Cauchy
comparison may be chosen separately, since the dilation coefficients
$c_{i,N}$ and $B_N$ are independent of translations of the blocking.

\emph{Approximation of continuous tests.}
Every microscopic descendant of a layer-$j$ cell is within
$CL^{-j}$ of its center.  Replace each continuous test by its value at
these cell centers.  Its error on a fixed compact set is bounded by
its modulus of continuity at $CL^{-j}$.  To estimate a mixed moment,
expand the difference of the two products by replacing one factor at
a time.  Nonnegativity of the component moment measures and
Equation~\eqref{eq:unit-box-majorant} bound each resulting term by that
modulus of continuity times a constant independent of $N$.  First let
$N\to\infty$ with $j$ fixed, and then let $j\to\infty$.  This proves
convergence for continuous tests.  Finite sums of product tests are
uniformly dense in the continuous functions on a product of compact
sets, and the total moment mass there is uniformly bounded.  Hence the
moment measures converge against every continuous compact test in the
insertion coordinates.

The factorial moment bounds give tightness of each finite list of
smeared real fields.  They also give uniform integrability of every
fixed polynomial in that list.  Any subsequential law consequently
has the moments just obtained and an exponential moment in a
neighborhood of zero.  Such a law is determined by its moments: its
moment generating function is analytic in that neighborhood, and its
Taylor coefficients are those moments.  Thus all subsequential laws
coincide and the finite-dimensional laws converge.

\begin{figure}[tb]
\centering
\begin{tikzpicture}[scale=0.9,>=Latex,every node/.style={font=\small}]
\begin{scope}[shift={(-4.8,0)}]
  \draw[step=.4,gray!35,very thin] (-1.25,-1.25) grid (1.25,1.25);
  \draw[->,thick] (0,0)--(1.6,0) node[right] {$e_1$};
  \draw[->,thick] (0,0)--(0,1.6) node[above] {$e_2$};
  \node[align=center] at (0,-1.9) {First regulator\\orientation $0$};
\end{scope}
\begin{scope}[rotate=53.130102]
  \draw[step=.65,MidnightBlue!65,thin] (-1.3,-1.3) grid (1.3,1.3);
  \draw[->,MidnightBlue,very thick] (0,0)--(1.6,0);
  \draw[->,MidnightBlue,very thick] (0,0)--(0,1.6);
\end{scope}
\node[align=center,text=MidnightBlue] at (0,-2.2)
  {Common coarse frame\\orientation $\theta$};
\begin{scope}[shift={(4.8,0)},rotate=106.260204]
  \draw[step=.4,gray!35,very thin] (-1.25,-1.25) grid (1.25,1.25);
  \draw[->,thick] (0,0)--(1.6,0);
  \draw[->,thick] (0,0)--(0,1.6);
\end{scope}
\node[align=center] at (4.8,-1.9) {Second regulator\\orientation $2\theta$};
\draw[->,thick] (-2.9,.25)--(-1.85,.25)
  node[midway,above=5pt,align=center] {$+\theta$};
\draw[->,thick] (2.9,.25)--(1.85,.25)
  node[midway,above=5pt,align=center] {$-\theta$};
\node at (0,2.15) {$\cos\theta=3/5,\qquad\sin\theta=4/5$};
\end{tikzpicture}
\caption{The two inclined first blockings produce the same physical
frame. Their relative inclinations are reflections of one another, so
their scalar field normalizations agree. After the first step, common
regular blockings erase the difference of their remaining data. The
drawing shows orientations only; the first coarse spacing is $5L$ times
the microscopic spacing.}
\label{fig:orientations}
\end{figure}


\emph{The two inclined blockings.}
Consider the fixed tilted square torus of side $5M$ and axes $O_+$
from Section~\ref{sec:trace}.  Put one microscopic regulator in the
standard axes and the other in axes $O_+^2$.  In their respective axes
use the inclined first steps $O_+$ and $O_-=O_+^{-1}$.  Their frames
then agree, because $O_+^2O_-=O_+$;
Figure~\ref{fig:orientations} shows these orientations.  Both regulators have the same
$\beta_N,a_N,B_N$ and compatible periods.

Stop at a fixed layer $j<N$.  If $B_N^*$ and $c_{i,N}^*$ denote the
inclined source normalizations, the exact source multiplier is now
\begin{equation}
       25L^{-2j}\frac{B_N}{B_N^*}
                    \prod_{i=1}^j(c_{i,N}^*)^{-1}.
 \label{eq:inclined-cell-source-multiplier}
\end{equation}
The factor $25$ comes from the area of the first block: the product of
the area factors through layer $j$ is $25L^{2(N-j)}$.  Reflection
covariance of the two inclined prescriptions makes the multipliers
in Equation~\eqref{eq:inclined-cell-source-multiplier} identical.
They are bounded at each fixed $j$ by
Proposition~\ref{prop:source-products} and the bounds on the individual
$c_{i,N}^*$.  Convergence of the ratio $B_N/B_N^*$ is not required.

Apply Equation~\eqref{eq:orig-14} and
Corollary~\ref{cor:source-endpoints} after the first step.  At the fixed
bottom layer, their generating functions differ by a quantity tending
to zero uniformly near the source origin, also when evaluated at the
common bounded multiplier
\eqref{eq:inclined-cell-source-multiplier}.  Explicitly, if $r_j>0$
is a common coarse-source comparison radius and $M_j$ bounds that
multiplier, restrict the physical cell sources to radius
$r_j/(2\max\{1,M_j\})$.  Their substituted arguments then lie in
radius $r_j/2$ for every $N$, so uniform convergence and Cauchy's
formula apply on one fixed polydisc without convergence of the
multiplier itself.  The common cell centers can
be chosen identical by translating the microscopic lattices.  Their
physical mesh can be fixed independently of $N$ at the chosen layer.
Every assigned microscopic point lies in its inclined first cell, so
the final descendants are within $C\,5L^{-j}$ of the common centers.
The same continuous-test approximation therefore proves equality of
the limiting moments, and hence of the finite-dimensional laws, on
this fixed tilted torus.  This step uses only the moment bounds for a
fixed tilted size.

\emph{Exchange of the cutoff and volume limits.}
For real $t_1,\ldots,t_r$ and compactly supported real tests
$f_1,\ldots,f_r$, set
\[
       F_N=\exp\!\left(i\sum_{\nu=1}^rt_\nu\phi_N(f_\nu)\right).
\]
It is a bounded local microscopic observable with $|F_N|=1$,
regardless of the size of $B_N$.  Once its support is contained in the
prescribed interior portion of an axis torus at doubling index $k$,
Section~\ref{sec:trace} gives
\begin{equation}
 \left|\E_{N,\mathrm{plane}}F_N
          -\E_{N,\mathrm{torus}(k)}F_N\right|
       \le C e^{-c2^k},
 \label{eq:characteristic-volume-comparison}
\end{equation}
uniformly in $N\ge K_0$.  The fixed-torus continuum limit has already
been proved.  The uniform moment bounds make the plane laws of this finite list tight.
For any subsequential plane limit, taking $N\to\infty$ in
Equation~\eqref{eq:characteristic-volume-comparison} bounds its
characteristic function within $Ce^{-c2^k}$ of the same fixed-torus
characteristic function.  Sending $k\to\infty$ identifies every
subsequential plane law.  This proves plane convergence and identifies
it with the successive torus and cutoff limit.  The same inequality
applied with $k=k(N)\to\infty$ gives every simultaneous axis exhaustion
claimed in the statement.  Uniform factorial moments give uniform
integrability, so moments converge along these limits as well.

For either regulator orientation, the trace-with-shift estimate of
Section~\ref{sec:trace} gives
Equation~\eqref{eq:characteristic-volume-comparison} also for the tilted
tori.  First take their fixed-size continuum limit, where the two
orientations have just been compared, and then send the doubling index
to infinity.  The resulting plane characteristic functions agree.
Because the comparison uses bounded characteristic functions, it
requires no moment bound uniform in the increasing tilted physical
size.
\end{proof}

\subsection{Euclidean symmetry and regularity}

Write $S_n^{\boldsymbol\alpha}$ for the plane moment measures just
constructed, and $S_n$ for the finite collection of their components.
There is no pointwise prescription for these correlation functions;
the following bounds specify their distributional meaning.

\begin{proposition}[Euclidean regularity]
\label{prop:euclidean-regularity}
The plane correlation functions are nonnegative Radon measures
componentwise, and for every scalar Schwartz test
$F\in\mathcal S((\R^2)^n)$ they satisfy
\begin{equation}
 |S_n^{\boldsymbol\alpha}(F)|
   \le C^n n!\,
       \sup_{x\in(\R^2)^n}
         \bigl[(1+|x|^2)^{2n}|F(x)|\bigr].
 \label{eq:orig-35}
\end{equation}
The lattice moment measures converge on Schwartz tests, with the same
bound.  Finite component norms can be incorporated by changing $C$.
The limiting correlations have permutation symmetry, internal $O(3)$
covariance, and invariance under all Euclidean isometries of $\R^2$.
Bounds for products of separately translated tests are uniform in
their separate translation parameters when the decay weights are
centered at those parameters.
\end{proposition}

\begin{proof}
The compact-test convergence and positivity from
Proposition~\ref{prop:continuum-limit} give nonnegative locally finite
measures.  Partition $\R^2$ into unit boxes indexed by $u\in\Z^2$.
On each such box the weights $(1+|x|^2)^{-2}$ and
$(1+|u|^2)^{-2}$ are comparable by a fixed constant.  Therefore
Equation~\eqref{eq:unit-box-majorant}, followed by summation of these
weights, gives
\[
 |S_n^{\boldsymbol\alpha}(F)|
  \le C^n n!\,
       \sup_{x_1,\ldots,x_n}
          \left[\prod_{i=1}^n(1+|x_i|^2)^2
                          |F(x_1,\ldots,x_n)|\right].
\]
Here the sum over product boxes factors into $n$ convergent sums, so
its cost is exponential in $n$.  Since
$\prod_i(1+|x_i|^2)^2\le(1+\sum_i|x_i|^2)^{2n}$, this proves
Equation~\eqref{eq:orig-35}.  The proof applies uniformly to the lattice
measures.  Cutting off a Schwartz test outside a large compact set
then gives convergence on Schwartz space: the weighted supremum of
the discarded tail tends to zero.  Replacing each weight by
$(1+|x_i-y_i|^2)^{-2}$ proves the assertion concerning separate
translations, with constants independent of $y_1,\ldots,y_n$.

Permutation symmetry and internal $O(3)$ covariance pass directly
from the lattice.  For a fixed physical translation, approximate its
displacement by lattice displacements.  The error on a smooth test
tends to zero by the majorants and test-space continuity.  Thus
translation invariance also passes to the limit.  The microscopic
reflections and quarter turns similarly give reflection and square
symmetry.

Let $\theta$ be the angle of $O_+$, so
$\cos\theta=3/5$.  Rotating the entire regulator rotates the arguments
of its moment functions.  The last assertion of
Proposition~\ref{prop:continuum-limit} therefore gives invariance under
rotation by $2\theta$.  This angle generates a dense subgroup of the
circle.  Indeed, if $\theta/\pi$ were rational, then
$e^{i\theta}$ would be a root of unity and
$2\cos\theta=6/5$ would be a rational algebraic integer, hence an
integer, a contradiction.  The same irrationality holds for
$2\theta/(2\pi)$.  Rotations act continuously on Schwartz test space;
Equation~\eqref{eq:orig-35} consequently extends this dense-subgroup
invariance to all rotations.  Together with reflections and
translations this is Euclidean invariance.
\end{proof}

The fourth-cumulant test in Section~\ref{sec:os} uses cell indicators before
approximating them by smooth functions.  The next estimate justifies
that approximation and also allows the microscopic cell boundaries to
move slightly with $N$.

\begin{lemma}[Boundary strips]
\label{lem:boundary-strips}
Fix an order $n$, a bounded region $K\subset\R^2$, and a bounded axis
rectangle $A$.  Let
$E_\varepsilon=K\cap\{x:\operatorname{dist}(x,\partial A)<\varepsilon\}$.
In the plane, uniformly in $N\ge K_0$ and in the microscopic origin,
\begin{equation}
 S_{n,N}^{\boldsymbol\alpha}
       (E_\varepsilon\times K^{n-1})
       \le C_{n,K,A}(\varepsilon+a_N).
 \label{eq:boundary-strip-bound}
\end{equation}
The same estimate holds on the axis tori, with the sets interpreted in
fixed bounded periodic charts.  It also holds with any one insertion
in the strip.  Consequently continuum convergence extends to products
of bounded axis-box indicators, to smooth approximations of these
indicators, and to cell boxes whose boundaries converge to the given
box boundaries.  More generally the same conclusion holds for bounded
Jordan measurable sets.
\end{lemma}

\begin{proof}
Choose a fixed bounded relative-coordinate box $D$ containing $K-K$.
Use coordinates in the regulator axes, so the full displacement lattice
is $a_N\Z^2$, regardless of the microscopic origin.  On a torus evaluate
translated spins periodically.  For a lattice site $x$, define
\[
 A_N(x)=(B_Na_N^2)^n
     \sum_{y_2,\ldots,y_n\in a_N\Z^2\cap D}
       \E\!\left[q_x^{\alpha_1}
                    \prod_{i=2}^n q_{x+y_i}^{\alpha_i}\right].
\]
For $n=1$ the sum and product over the other insertions are empty.
Every term is nonnegative.  Microscopic translation invariance makes
$A_N(x)$ independent of $x$; denote its value by $A_N$.  Notice that
all $n$ field-normalization factors and all $n$ lattice area factors
are included in $A_N$.

Let $Q$ be a fixed unit box.  Summing the first insertion over its
lattice sites and using positivity gives
\[
  \#(Q\cap\Lambda_N) A_N
     \le S_{n,N}^{\boldsymbol\alpha}
                  (Q\times(Q+D)^{n-1})
     \le C_{n,K}.
\]
The last inequality follows by a fixed unit-box covering and
Equation~\eqref{eq:unit-box-majorant}.  Since
$\#(Q\cap\Lambda_N)\ge c a_N^{-2}$, with a constant uniform over the
lattice origin, we obtain $A_N\le C_{n,K}a_N^2$.  If the first insertion
is in $E\subset K$, every choice of the other insertions in $K$ is
included in this relative-coordinate sum.  Hence
\begin{equation}
 S_{n,N}^{\boldsymbol\alpha}(E\times K^{n-1})
       \le C_{n,K}a_N^2\#(E\cap\Lambda_N).
 \label{eq:first-insertion-count}
\end{equation}
For a fixed rectangle boundary, elementary lattice counting gives
$a_N^2\#(E_\varepsilon\cap\Lambda_N)
 \le C_{K,A}(\varepsilon+a_N)$, proving
Equation~\eqref{eq:boundary-strip-bound}.

On a torus use the same displacement box $D$, with periodic evaluation
of $x+y_i$.  A fixed bounded displacement box can represent a periodic
site more than once, but its multiplicity is bounded at the fixed
physical scale under consideration.  The unit-box argument and the
moment bound therefore give the same estimate, with this fixed
multiplicity included in the constant.  Permutation symmetry treats
any insertion coordinate.

For completeness, if $A$ is merely bounded and Jordan measurable,
the lattice squares centered at points of
$\{\operatorname{dist}(x,\partial A)<\varepsilon\}$ are contained in the
$(\varepsilon+Ca_N)$-neighborhood of $\partial A$.  Thus the counting
term in Equation~\eqref{eq:first-insertion-count} is bounded by a
constant times the area of that neighborhood, which tends to zero as
$\varepsilon\downarrow0$ and $N\to\infty$.  Continuous inner and outer
approximations to the indicator now have arbitrarily small error in
every fixed-order correlation.  The error for a product of indicators
is bounded by the sum of the errors with one argument in a boundary
neighborhood.  Compact-test convergence applies to the continuous
approximants, and then the neighborhood widths tend to zero.  The same
argument covers moving boxes, since their symmetric differences lie
in shrinking neighborhoods of the limiting boundaries.
\end{proof}
